
% --- Aprendizaje automático y federado ---
\newacronym{sgd}{SGD}{Stochastic Gradient Descent}
\newacronym{fl}{FL}{Federated Learning}
\newacronym{fedavg}{FedAvg}{Federated Averaging}
\newacronym{fedprox}{FedProx}{Federated Proximal}
\newacronym{fedopt}{FedOpt}{Federated Optimization}
\newacronym{iid}{IID}{Independent and Identically Distributed}

% --- Hardware y cómputo ---
\newacronym{cpu}{CPU}{Central Processing Unit}
\newacronym{gpu}{GPU}{Graphics Processing Unit}
\newacronym{ram}{RAM}{Random Access Memory}
\newacronym{vram}{VRAM}{Video Random Access Memory}
\newacronym{cuda}{CUDA}{Compute Unified Device Architecture}
\newacronym{mpi}{MPI}{Message Passing Interface}
\newacronym{nvml}{NVML}{NVIDIA Management Library}

% --- Software, comunicación y formatos ---
\newacronym{api}{API}{Application Programming Interface}
\newacronym{sdk}{SDK}{Software Development Kit}
\newacronym{grpc}{gRPC}{gRPC Remote Procedure Call}
\newacronym{tls}{TLS}{Transport Layer Security}
\newacronym{json}{JSON}{JavaScript Object Notation}

% --- Seguridad y privacidad ---
\newacronym{dp}{DP}{Differential Privacy}
\newacronym{psi}{PSI}{Private Set Intersection}

% --- Normativa y organismos ---
\newacronym{iso}{ISO}{International Organization for Standardization}
\newacronym{w3c}{W3C}{World Wide Web Consortium}

% --- Académicos / del trabajo ---
\newacronym{oe}{OE}{Objetivo Específico}
\newacronym{tfg}{TFG}{Trabajo de Fin de Grado}

% -----------------------------------------------------------------------------
% GLOSARIO DE TÉRMINOS
% Definiciones propias y concisas de los conceptos empleados en la memoria.
% Cada entrada incluye 'sort' en ASCII para evitar errores de makeindex.
% -----------------------------------------------------------------------------

% --- Conceptos de aprendizaje federado ---
\newglossaryentry{aprendizaje-federado}{
  name={aprendizaje federado},
  sort={aprendizaje federado},
  description={Paradigma de aprendizaje automático que entrena un modelo de forma colaborativa entre varios clientes sin centralizar sus datos; en lugar de los datos, solo se intercambian las actualizaciones del modelo}
}

\newglossaryentry{fl-horizontal}{
  name={aprendizaje federado horizontal},
  sort={aprendizaje federado horizontal},
  description={Variante en la que los clientes comparten el mismo espacio de características pero poseen ejemplos distintos. Es el escenario empleado en este trabajo}
}

\newglossaryentry{fl-vertical}{
  name={aprendizaje federado vertical},
  sort={aprendizaje federado vertical},
  description={Variante en la que los clientes comparten parte de los individuos o entidades pero describen atributos distintos, lo que exige alinear registros}
}

\newglossaryentry{cross-silo}{
  name={cross-silo},
  sort={cross-silo},
  description={Escenario federado con pocos clientes estables y con capacidad de cómputo elevada (por ejemplo hospitales o bancos), con participación habitualmente completa}
}

\newglossaryentry{cross-device}{
  name={cross-device},
  sort={cross-device},
  description={Escenario federado con un número muy elevado de dispositivos de usuario final, con disponibilidad intermitente, recursos limitados y conexiones inestables}
}

\newglossaryentry{ronda-federada}{
  name={ronda federada},
  sort={ronda federada},
  description={Iteración completa del ciclo de distribución del modelo global, entrenamiento local en los clientes, envío de los modelos locales y agregación en el servidor}
}

\newglossaryentry{epoca-local}{
  name={época local},
  sort={epoca local},
  description={Número de veces que cada cliente recorre su conjunto de datos local durante el entrenamiento de una ronda}
}

\newglossaryentry{tamano-lote}{
  name={tamaño de lote},
  sort={tamano de lote},
  description={Número de ejemplos que se procesan antes de cada actualización de los parámetros del modelo (\textit{batch size})}
}

\newglossaryentry{tasa-aprendizaje}{
  name={tasa de aprendizaje},
  sort={tasa de aprendizaje},
  description={Hiperparámetro que regula la magnitud de cada paso de actualización de los parámetros del modelo (\textit{learning rate})}
}

\newglossaryentry{estrategia-agregacion}{
  name={estrategia de agregación},
  sort={estrategia de agregacion},
  description={Procedimiento del servidor para combinar los modelos locales recibidos y obtener una nueva versión del modelo global}
}

\newglossaryentry{datos-no-iid}{
  name={datos no IID},
  sort={datos no IID},
  description={Datos no independientes ni idénticamente distribuidos; situación en la que la distribución de los datos difiere de un cliente a otro, lo que dificulta la convergencia del modelo global}
}

% --- Algoritmos federados ---
\newglossaryentry{alg-fedavg}{
  name={FedAvg (algoritmo)},
  sort={FedAvg},
  description={Algoritmo de referencia del aprendizaje federado que agrega los modelos locales mediante un promedio ponderado por el número de ejemplos de cada cliente}
}

\newglossaryentry{alg-fedprox}{
  name={FedProx (algoritmo)},
  sort={FedProx},
  description={Variante de FedAvg que añade un término proximal al entrenamiento local para penalizar la desviación de los pesos locales respecto al modelo global}
}

\newglossaryentry{alg-fedopt}{
  name={FedOpt (familia)},
  sort={FedOpt},
  description={Familia de optimización federada adaptativa que sustituye el promedio del servidor por un optimizador con momento y tasa de aprendizaje propios}
}

\newglossaryentry{alg-fedadam}{
  name={FedAdam},
  sort={FedAdam},
  description={Variante de FedOpt que emplea el optimizador Adam en el lado del servidor}
}

\newglossaryentry{alg-fedadagrad}{
  name={FedAdagrad},
  sort={FedAdagrad},
  description={Variante de FedOpt que emplea el optimizador Adagrad en el lado del servidor}
}

\newglossaryentry{alg-scaffold}{
  name={SCAFFOLD},
  sort={SCAFFOLD},
  description={Algoritmo federado que corrige la deriva entre clientes mediante variables de control para mejorar la convergencia con datos heterogéneos}
}

% --- Modelos, datos y entrenamiento ---
\newglossaryentry{cifar10}{
  name={CIFAR-10},
  sort={CIFAR-10},
  description={Conjunto de 60\,000 imágenes en color de 32 por 32 píxeles distribuidas en diez clases, empleado en este trabajo como banco de pruebas de clasificación}
}

\newglossaryentry{resnet18}{
  name={ResNet-18},
  sort={ResNet-18},
  description={Red neuronal residual de 18 capas; en este trabajo se adapta a las dimensiones de CIFAR-10 sustituyendo la primera convolución y eliminando el \textit{maxpool} inicial}
}

\newglossaryentry{funcion-perdida}{
  name={función de pérdida},
  sort={funcion de perdida},
  description={Función que cuantifica el error del modelo comparando su predicción con la salida esperada; en este trabajo se emplea la entropía cruzada}
}

% --- Seguridad y privacidad ---
\newglossaryentry{privacidad-diferencial}{
  name={privacidad diferencial},
  sort={privacidad diferencial},
  description={Técnica que añade ruido controlado a los datos o a las actualizaciones del modelo para limitar la información que puede inferirse sobre un individuo concreto}
}

\newglossaryentry{agregacion-segura}{
  name={agregación segura},
  sort={agregacion segura},
  description={Protocolo criptográfico que permite combinar las actualizaciones de los clientes sin que el servidor observe cada una por separado}
}

\newglossaryentry{cifrado-homomorfico}{
  name={cifrado homomórfico},
  sort={cifrado homomorfico},
  description={Esquema de cifrado que permite operar sobre datos cifrados sin necesidad de descifrarlos previamente}
}

% --- Infraestructura y ejecución ---
\newglossaryentry{framework}{
  name={framework},
  sort={framework},
  description={Conjunto de componentes reutilizables que proporciona la infraestructura para construir y ejecutar un tipo de sistema, en este caso de aprendizaje federado}
}

\newglossaryentry{backend}{
  name={backend},
  sort={backend},
  description={Motor de ejecución subyacente de un framework; en FedML determina si el entrenamiento es secuencial (\textit{sp}) o concurrente (MPI)}
}

\newglossaryentry{ray}{
  name={Ray},
  sort={Ray},
  description={Framework de cómputo distribuido sobre el que Flower ejecuta su motor de simulación de clientes}
}

\newglossaryentry{modo-secuencial}{
  name={modo secuencial},
  sort={modo secuencial},
  description={Modo de ejecución en el que se mantiene una sola tarea de entrenamiento activa cada vez, de manera que los clientes de una ronda se procesan uno tras otro}
}

\newglossaryentry{modo-concurrente}{
  name={modo concurrente},
  sort={modo concurrente},
  description={Modo de ejecución en el que varios clientes o procesos permanecen activos simultáneamente, lo que permite observar el paralelismo y su impacto en los recursos}
}

% --- Ejes de evaluación ---
\newglossaryentry{usabilidad}{
  name={usabilidad},
  sort={usabilidad},
  description={Eje de evaluación que valora la experiencia de desarrollo de cada framework: instalación, configuración, claridad de los registros y facilidad de depuración}
}

\newglossaryentry{flexibilidad}{
  name={flexibilidad},
  sort={flexibilidad},
  description={Propiedad arquitectónica que describe hasta qué punto un framework permite adaptarse a escenarios distintos de los previstos por defecto, modificando, sustituyendo o ampliando su comportamiento}
}
